\section{Exactness and the rational monodromy equation}

% BEGIN K85F1
\paragraph{K85F1. Principal parts and completeness on the common tame field.}
Use the integral coefficient ring $O$ of Q84T1 after extension, with
fraction field $L$. Write $A=O[[u]]$, $E=u^e+p$ and
$M$ for the original free module over the product of the $a$ copies
of $A$. Write $\mathcal L$ for its Frobenius-column lattice.
Let $R_\pi$ be the ring of analytic functions on the closed disc
$|u|\le p^{-1/e}$. The principal-parts map is
$y\mapsto[y]$ modulo $M_{\rm an}=R_\pi M$. On $E^{-b}M$
its kernel is $M$. Indeed divide each numerator by the distinguished
monic $E^b$ over $O[[u]]$. If the fraction belongs to $R_\pi$,
the polynomial remainder has zeros of multiplicity at least $b$
at every root of $E$, after scalar extension. The roots are distinct
and on the disc. Its degree is less than $eb$, so it vanishes.
Over $R_\pi$, division by the same polynomial gives the same
space of remainders over $L$ (equivalently rescale to the closed
unit disc and use distinguished division). Hence
\[
 Z=\big(E^{-b}M/M\big)^\Delta,\qquad
 Z_L=\big(E^{-b}M_{\rm an}/M_{\rm an}\big)^\Delta=Z[1/p].
\]
For the second equality multiply the finitely many coefficients
of a remainder by a common power of $p$. Invariance of its class
survives this multiplication. The unprojected integral quotient
is free with basis $E^{-b}u^j$ in each coordinate for $0\le j<eb$.
The invariant submodule is finite torsion free over the DVR $O$,
hence free and $p$-adically separated. To check completeness of
$Y'=u^e(E^{-b}M)^\Delta$, identify it via multiplication by
$E^b/u^e$ with $M^\Delta$ (the scalar is tame invariant).
The coefficients of a compatible family modulo $p^n$ in $M$ have
unique limits coefficient by coefficient, which assemble to a
power series over $O$. The group action preserves $p^nM$ and
therefore fixes the limit of invariant approximants. For every $n$
one has $p^nM\cap M^\Delta=p^n M^\Delta$ by torsion freeness.
This proves completeness in the required module topology.
% END K85F1

% BEGIN K85F2
\paragraph{K85F2. Exactness on the integral principal parts.}
Let $M_0\xrightarrow{F}M_1\xrightarrow{G}M_2$ be an exact
triple in the sense of Q84T3. Suppose $p^s M_2\subset G M_1$.
Choose a linear $J:M_2\to M_1$ lifting $p^s$ on a basis.
As $G(p^s-JG)=0$, exactness and injectivity of $F$ give a linear
$B:M_1\to M_0$ with $FB=p^s-JG$. Applying this to $F$ gives
$BF=p^s$. In particular the induced map from all principal parts is
injective: if $Fx$ is analytic, $x=p^{-s}BFx$ is analytic.
At $u=0$ the same relation gives injectivity of each component;
since $F$ intertwines the inertia actions there, each kernel character
occurs in the middle module.

Take $q=[y]\in Z_1$ with $Gq=0$ and $y\in E^{-b}M_1$.
Then $Gy\in E^{-b}M_2\cap M_{2,\rm an}=M_2$ by K85F1.
Pick $x\in M_1$ with $Gx=p^sGy$. There is $m\in M_0$
with $Fm=E^b(p^s y-x)$, by exactness. Thus $[E^{-b}m]$ maps
to $p^s q$. This class is tame fixed: the image of its difference
with any conjugate vanishes, and $F$ is injective on principal
parts. We have lifted the entire integral obstruction at the same
cost $p^s$.
% END K85F2

% BEGIN K85F3
\paragraph{K85F3. Strictness at each filtration jump.}
Write $M_A=M[1/p]$, $L_A=\mathcal L[1/p]$ and use a subscript
$\pi$ for extension to $R_\pi$. The filtration is the preimage
in $L_A$ of the image of $L_A\cap E^i M_A$ modulo $EL_A$.
It can be tested analytically. For $x\in L_\pi\cap E^iM_\pi$,
divide each coordinate of $E^{-i}x$ by a sufficiently high power
$E^m$, obtaining a polynomial remainder $g$. Choose $m$ so that
$E^{i+m}M_\pi\subset EL_\pi$, using the two height bounds.
Then $x-E^ig\in EL_\pi$ and $E^ig\in L_\pi$. In Frobenius
coordinates the entries of $E^ig$ have only powers of $E$ as
denominators over $A[1/p]$. Analyticity cancels those powers by
division as in K85F1, clearing one coefficient denominator first.
Thus $E^ig\in L_A\cap E^iM_A$. The same division shows
$L_A\cap EL_\pi=EL_A$.

Set $D=E\nabla$, where $\nabla$ is the Euler connection with scalar
derivation $u\,d/du$. On the analytic Frobenius lattice the ordinary
connection preserves $L_\pi$: on its generating columns its values
lie in the divided-power extension, whose coefficients converge
on this disc. Since $uE'=e(E-p)$, the induced action of $D$ on
$E^jL_\pi/E^{j+1}L_\pi$ is $-jep$. On the physical lattice,
$DM_\pi\subset M_\pi$ is the analytic stability in the
Breuil--Kisin datum. Leibniz then preserves $E^iM_\pi$ for every
integer $i$. Choose $m=\max(0,i+b)$ and put
\[
 P_m(X)=\prod_{j=1}^m(X+jep),\qquad P_m(0)=(ep)^m m!\ne0 .
\]
Suppose $x\in L_{0,A}$ and $Fx$ belongs to filtration $i$.
Decompose $Fx$ as an element of $E^iM_{1,\pi}$ plus an element of
$EL_{1,\pi}$. On the latter the successive factors of $P_m(D)$
raise the $E$-order to at least $m+1$. By $E^bL_{1,\pi}\subset
M_{1,\pi}$ this term too is now in $E^iM_{1,\pi}$. By
horizontality and the matrix $B$ in K85F2 we get
$P_m(D)x\in E^iM_{0,\pi}$. It lies in $L_{0,\pi}$ and reduces to
$P_m(0)x$ modulo $EL_{0,\pi}$. The analytic filtration test and
division by the nonzero scalar give membership of $x$ in filtration
$i$ of $L_{0,A}$. Applying this reflection at $i+1$, and forward
functoriality at $i$, shows that equality of the two adjacent
filtration steps in the middle forces equality in the kernel.
In particular all kernel jumps belong to the middle jump set.
% END K85F3

% BEGIN K85F4
\paragraph{K85F4. A rational primitive and the class in the Euler image.}
Write $\mathcal S$ for the completed divided-power coefficient ring
with factorial weights $\lfloor n/e\rfloor !$, and
$\mathcal M_{\rm dp}=\mathcal S\mathcal L$. For a Frobenius
module with matrix $A(u)$ let $\xi$ denote its horizontal section
at $u=0$. In Frobenius coordinates $\xi$ and the inverse change
of basis have entries in $\mathcal S[1/p]$. Here are the
convergence details. Choose a radius $r=p^{-k}$ with
$1/(ep)<k<1/(e(p-1))$. In the normalized Frobenius product the
initial factor $A(u)$ is removed by passing to Frobenius coordinates;
the later factors $A(u^{p^j})$ and their inverses are analytic on
the disc of radius $r$ for $j\ge1$, since $r^p<p^{-1/e}$.
With the component indices cycled, each new normalized factor differs
from its constant value by $O(r^{p^j})$ in Gauss norm, up to a
fixed exponential $C^j$ from conjugating by constant matrices.
This follows by factoring $u$ from the difference at zero on a
fixed disc containing all the arguments $u^{p^j}$ and adjusting
the constant for its radius. The normalized products and the inverse
products therefore converge in the Gauss algebra and are mutual
inverses. Their formal limits give the section by the Frobenius
recursion and its normalization at zero. The ring on that disc
embeds in $\mathcal S[1/p]$ since
\[
 |\lfloor n/e\rfloor !|_p
 \le C(n+1)p^{-n/(e(p-1))}.
\]
This follows by the base-$p$ digit-sum formula for the valuation
of a factorial. Thus the divided-power weighted coefficients tend
to zero and have a common denominator.

First take $y\in Y$ and form the triangular Frobenius extension
with extra generator $x$, whose Frobenius coordinate is $(y,1)$
in the convention of the extension matrix. In the notation of
that extension put $t=\varphi x-\xi(\operatorname{red}\varphi x)$.
It belongs to the enlarged divided-power lattice after inverting
$p$ by the preceding calculation. It has zero reduction at zero.
Its Euler derivative is the inclusion of $\zeta(y)$ by
differentiating the defining Frobenius recursion. The projection of
$t$ to the trivial line is therefore horizontal with zero reduction
and vanishes. Indeed in horizontal coordinates the Euler equation
is $u f'=0$, whose solutions in characteristic zero are constants.
Consequently $t$ is the image of a vector $s\in
\mathcal M_{\rm dp}[1/p]$ of the submodule: the last Frobenius
column has last coordinate one, forcing its coefficient to be zero.
The tame action fixes $s$ since it fixes $x$, commutes with
Frobenius and transports the horizontal section. We have
$\nabla s=\zeta(y)$. For $y\in Y[1/p]$ clear a scalar
denominator and divide it out after this construction.

The class to use is $[y-s]$. Both $E^b y$ and $E^b s$ belong to
$M_{\rm an}$ by the height bound and the analytic inclusion
$\mathcal S[1/p]\subset R_\pi$ on this smaller disc. Thus
$[y-s]\in Z_L$. By the definitions
\[
 L_c[y-s]=[E\nabla y-E\zeta(y)]=\mathrm{Mon}(y).
\]
This supplies the rational image assertion even when $s$ itself
does not give an integral principal part.
% END K85F4
